% Part: applied-modal-logics
% Chapter: epistemic-logic
% Section: bisimulations
%READERNOTE{SOL6-A190: द्विअनुकरणासाठी graphs समरूप असण्याची गरज नाही; forth/back उत्तरवर्ती अनुरूपता आवश्यक आहे हे स्पष्ट केले. चित्रात न दिलेल्या valuations आणि इतर कर्त्यांचे संबंध नेमके निवडल्यावरच दाखवलेला तीन-जोड्यांचा संबंध द्विअनुकरण आहे असे मांडले; मूळ graph जपला.}

\documentclass[../../../include/open-logic-section]{subfiles}

\begin{document}

\olfileid{aml}{el}{bsd}

\olsection{द्विअनुकरणे}

आपल्या ज्ञानविषयक भाषेच्या अभिव्यक्तिक्षमतेविषयी
उरलेला एक प्रश्न प्रतिमाने आणि त्यांत सत्य असलेली
सूत्रे यांमधील संबंधाशी निगडित आहे. चौकट-अनुरूपतेच्या
निष्कर्षांवरून आपण पाहिले की काही सूत्रे एखाद्या
चौकटीत वैध असतील, तर त्या चौकटीला विशिष्ट गुणधर्म
असतात. पण उदाहरणार्थ, एखाद्या जगापासून त्याच
जगाकडे जाणारा परावर्ती बाण आणि प्रत्येक जग पुढच्या
जगाकडे जाणारी अनंत साखळी यांमध्ये भेद करता येईल
अशी आपली मोडल भाषा आहे का? म्हणजे, या दोन
जगांपैकी फक्त एकाच जगात सत्य असेल असे एखादे
सूत्र~$A$ आहे का?

संबंधी प्रतिमानांतील संबंधित जगांत मूलभूत
ज्ञानविषयक सूत्रांनी भेद करता येत नाही, ही कल्पना
द्विअनुकरण संबंध पकडतो. यासाठी दोन प्रतिमानांतील
जगांची संख्या किंवा त्यांचे ग्राफ समरूप असण्याची
अट नाही; पुढील अटी योग्य अनुरूपता सांगतात.

\begin{defn}[द्विअनुकरण]
$M_1 = \tuple{W_1, R_1, V_1}$ आणि $M_2 = \tuple{W_2, R_2, V_2}$
ही एकाच कर्तृचिन्ह-संच~$G$ व एकाच विधानीय
भाषेची दोन संबंधी प्रतिमाने असू द्या. तसेच
$\mathcal{R} \subseteq W_1 \times W_2$ हा द्विस्थानी
संबंध असू द्या. प्रत्येक $\tuple{w_1, w_2} \in
\mathcal{R}$ साठी पुढील अटी पूर्ण होत असतील, तर
$\mathcal{R}$ हे \emph{द्विअनुकरण} आहे असे म्हणतो:

\begin{enumerate}
  \item सर्व !!{propositional variable}s~$p$ साठी
  $w_1 \in V_1(p)$ तेव्हा आणि केवळ तेव्हाच, जेव्हा
  $w_2 \in V_2(p)$.
  \item सर्व कर्ते $a \in G$ आणि सर्व जगे $v_1 \in W_1$
  यांसाठी, $R_{1_a} w_1 v_1$ असेल, तर $R_{2_a} w_2 v_2$
  आणि $\tuple{v_1, v_2} \in \mathcal{R}$ असलेले काही
  $v_2 \in W_2$ असते.
  \item सर्व कर्ते $a \in G$ आणि सर्व जगे $v_2 \in W_2$
  यांसाठी, $R_{2_a} w_2 v_2$ असेल, तर $R_{1_a} w_1 v_1$
  आणि $\tuple{v_1, v_2} \in \mathcal{R}$ असलेले काही
  $v_1 \in W_1$ असते.
\end{enumerate}

$M_1$ आणि $M_2$ यांमधील एखादे द्विअनुकरण जगे
$w_1$ आणि $w_2$ यांना जोडत असेल, तर आपण
$\tuple{M_1, w_1} \leftrightarroweq
\tuple{M_2, w_2}$ असेही लिहू शकतो. तेव्हा
$\tuple{M_1, w_1}$ आणि $\tuple{M_2, w_2}$ यांना
\emph{द्विअनुकरणसदृश} म्हणतो.
\end{defn}

द्विअनुकरण संबंधातील निरनिराळ्या अटी वेगवेगळ्या
गोष्टी सुनिश्चित करतात. पहिली अट सर्व !!{propositional variable}s
बाबतीत एकमत असल्याचे सुनिश्चित करते;
म्हणून द्विअनुकरणसदृश जगांत तीच मोडल-विरहित सूत्रे
सत्य असतात. उरलेल्या दोन अटींना अनुक्रमे ``पुढे''
आणि ``मागे'' अशा अटी म्हणतात. संबंधित जगांच्या प्रत्येक
प्राप्य उत्तरवर्ती जगाला दुसऱ्या प्रतिमानात संबंधित
प्राप्य उत्तरवर्ती जग दोन्ही दिशांनी मिळते; उत्तरवर्ती
जगांची संख्या सारखी किंवा बाणांच्या रचना समरूप असण्याची अट नाही.

\begin{thm}
$\tuple{M_1, w_1} \leftrightarroweq \tuple{M_2, w_2}$
असेल, तर प्रत्येक !!{formula}~$!A$ साठी
$\mSat{M_1}{!A}[w_1]$ तेव्हा आणि केवळ तेव्हाच,
जेव्हा $\mSat{M_2}{!A}[w_2]$.
\end{thm}

\begin{figure}
  \begin{center}
    \begin{tikzpicture}[modal]
      \node[world] (w1) {$w_1$};
      \node[world] (w2) [above left=of w1]{$w_2$};
      \node[world] (w3) [above right=of w1] {$w_3$};
      \draw[<->] (w1) to node {$a$} (w2);
      \draw[->,loop below] (w1) to node {$a$} (w1);
      \draw[<->] (w1) to [swap] node {$a$} (w3);
      \draw[->,loop above] (w2) to node {$a$} (w2) ;
      \draw[->,loop above] (w3) to node {$a$} (w3) ;

      \node[world](v1)[right of=w1, xshift=1.5in]{$v_1$};
      \node[world](v2)[above of=v1]{$v_2$};
      \draw[<->] (v1) to node {$a$} (v2);
      \draw[->,loop below] (v1) to node {$a$} (v1);
      \draw[->,loop above] (v2) to node {$a$} (v2) ;

      \draw[-,dotted] (w1) to (v1) ;
      \draw[-,dotted,bend left=45] (w2) to (v2) ;
      \draw[-,dotted,bend left=30] (w3) to (v2) ;
    \end{tikzpicture}
  \end{center}
  \caption{दोन द्विअनुकरणसदृश प्रतिमाने.}
  \ollabel{fig:bisimilar}
\end{figure}

\olref{fig:bisimilar} मध्ये बाणांची रचना दाखवली आहे;
द्विअनुकरणासाठी मूल्यांकनांची अनुरूपताही आवश्यक आहे.
येथे प्रत्येक !!{propositional variable}~$p$ साठी
$w_1$ आणि $v_1$ येथे त्याचे सत्यतामूल्य सारखे,
आणि $w_2,w_3,v_2$ या तिन्ही जगांत त्याचे सत्यतामूल्य
सारखे घेऊ. इतर कर्ते असतील, तर त्यांचे संबंध या
दोन प्रतिमानांत रिकामे घेऊ. मग
$\{\tuple{w_1,v_1},\tuple{w_2,v_2},\tuple{w_3,v_2}\}$
हा संबंध द्विअनुकरण आहे. म्हणून संबंधित जगांत मूलभूत
मोडल भाषेतील सूत्रांनी भेद करता येत नाही. मात्र
$w_2$ आणि~$w_3$ यांत निरनिराळी !!{propositional variable}s
सत्य असतील, तर ते दोन्ही~$v_2$ शी जोडणारा हाच
संबंध पहिली अट पूर्ण करत नाही.

\end{document}
